\clearpage
\SourceLeaf{338}
\pdfbookmark[2]{5. Gauss an Olbers, 7. Januar 1815}{parabolisch-005}
\begin{center}
[5.]\par
\medskip
\textsc{Gauss} an \textsc{Olbers}, Göttingen, 7. Januar 1815.
\end{center}
\noindent\hspace*{\parindent}.\ .\ .\ .\ .\ .\par
Ich habe Ihnen in meinem vorigen Briefe die Annäherungsformel für die Zwischenzeit
\[
t''-t=3mk\BareRoot{(r+r'')}\,.\BareRoot[10]{\left(1-\tfrac5{12}\frac{kk}{(r+r'')^2}\right)}\text{\begingroup\renewcommand{\thefootnote}{[*]}\footnotemark\endgroup}
\]
\HistoricalDisplayNoteText{Hier bezeichnet, abweichend von den vorstehenden Notizen, \(k\) die Chorde, \(\frac1{6m}\) die im Vorigen mit \(k\) bezeichnete Constante, vgl. S. 330.]}
geschrieben und gesagt, dass sie in den meisten Fällen hinreichend genau sei. Ich hätte sagen sollen, es werde in der Praxis kein Fall vorkommen, wo sie nicht überflüssig genau wäre. So lange \(k<\tfrac12(r+r'')\), afficirt der Fehler noch nicht die 5\textsuperscript{te} Decimale. Es ist eine Lust, zu sehen, wie bequem die Versuche für \(u\) danach durchgeführt werden. Jeder Versuch erfordert bloss 8 Aufschlagungen, 3 in den Logarithmen, 4 in den Sinustafeln und eine in meiner Hülfstafel für Logarithmen von Summen und Differenzen. Um so wenig wie möglich zu schreiben zu haben, setze ich \(r\BareRoot{\frac{12}5}=s\), \(r''\BareRoot{\frac{12}5}=s''\); dadurch wird, wenn man noch \(\frac{k}{s+s''}=\sin Q\) setzt:
\[
\frac{t''-t}{m\BareRoot[4]{\frac{135}4}}=\frac{k\BareRoot{(s+s'')}}{\BareRoot[5]{\sec Q}}\mathpunct{\cdot}
\]

Die ganze Rechnung besteht also vorher nur noch in folgendem:
\[
\begin{aligned}
\frac{bB}{u+c}&=\operatorname{tang}\theta,&\frac{b''B''}{u+c''}&=\operatorname{tang}\theta'',&\frac Au&=\operatorname{tang}\eta\\[5pt]
\frac{B\BareRoot{\frac{12}5}}{\sin\theta}&=s,&\frac{B''\BareRoot{\frac{12}5}}{\sin\theta''}&=s'',&\frac A{\sin\eta}&=\frac u{\cos\eta}=k.
\end{aligned}
\]

\SourceLeaf{339}\noindent\hspace*{\parindent}.\ .\ .\ .\ .\ .\par
Die ausserordentliche Leichtigkeit dieser Versuche macht, dass es ziemlich gleichgültig ist, mit was für einem Werthe von \(u\) man den Anfang macht. Indessen kann man doch auch hiebei einigermassen methodisch zu Werke gehen. Um nicht ganz im Blinden zu tappen, bestimme ich den ersten Werth von \(u\) auf folgende Art. Ich setze
\[
\frac{\BareRoot[6]{(A^4BB'')}}{\left(\frac{t''-t}{m\BareRoot{18}}\right)^{\frac23}}=\sin P,
\]
\[
u=-\tfrac16(c+c'')+\cos P\cdot\left(\frac{t''-t}{m\BareRoot{18}}\right)^{\frac23}\cdot\BareRoot[6]{(bb'')},
\]
was freilich nur eine höchst grobe Annäherung sein kann. Der Grund dieser Formel ist folgender. Ich abstrahire anfangs von \(\BareRoot[5]{\sec Q}\). Ich setze:
\[
\tfrac12(r+r'')=\BareRoot{\left\{\left(\frac{u+\frac12(c+c'')}{\BareRoot{(bb'')}}\right)^2+\BareRoot{(BB'')}^2\right\}}
\]
und, was freilich noch kühner ist,
\[
k^{\frac23}\left(\frac{r+r''}2\right)^{\frac13}=\BareRoot{\left\{\left(\frac{u+\frac16(c+c'')}{\BareRoot[3]{(\BareRoot{(bb'')})}}\right)^2+\left(A^{\frac23}(\BareRoot{(BB'')})^{\frac13}\right)^2\right\}}.
\]
Nachdem ich mit diesem ersten Werthe von \(u\) die Rechnung durchgeführt habe, erhalte ich den Fehler in Beziehung auf \(\log\frac{t''-t}{m\BareRoot[4]{\frac{135}4}}\mathpunct{\cdot}\) Jetzt mache ich die zweite Hypothese nicht auf gut Glück, sondern ich überlege, dass eine Veränderung von \(u\) bei mässigen Zwischenzeiten immer viel stärker auf \(\log k\) als auf \(\tfrac12\log(r+r'')\) wirkt, \(\sec Q\) gar nicht zu gedenken. Ich werfe also den ganzen Fehler der ersten Hypothese auf \(\log k\), und gehe so bei der zweiten Hypothese von dem geänderten Werthe von \(\log k\) aus, um danach \(u\) vermittelst der Formeln
\[
\frac Ak=\sin\eta,\qquad u=k\cos\eta
\]
zu bestimmen. Nachdem nun auch diese Hypothese durchgerechnet und ihr Fehler ausgemittelt ist, bestimme ich den dritten Werth von \(u\) schon durch Interpolation, und habe nun schon gewonnen Spiel. Sehen Sie hier die Anwendung auf das Beispiel meiner Abhandlung:

Die obige Formel gibt:
\[
\mathrm{I.}\quad u=0{,}15599.
\]

Aus der Zwischenzeit folgt \(\log\frac{t''-t}{m\BareRoot[4]{\frac{135}4}}=9{,}77932\), aus diesem Werthe \SourceLeaf{340}von \(u\) hingegen, wobei \(\log k=9{,}36029\) wird, folgt dafür \(9{,}65406\), also zu klein um 12526 Einheiten in der fünften Stelle. Ich nehme also für die zweite Hypothese \(\log k=9{,}48555\), woraus
\[
\mathrm{II.}\quad u=0{,}25561.
\]

Der Fehler dieser Hypothese wird \(+1598\), sie gibt nemlich \(9{,}79530\). Hienach finde ich durch Interpolation Verbesserung des zweiten \(u\): \(-0{,}01127\), also
\[
\mathrm{III.}\quad u=0{,}24434.
\]

Der Fehler dieser Hypothese ist \(+67\), also durch abermalige Interpolation zwischen II und III
\[
\mathrm{IV.}\quad u=0{,}24385,
\]
welcher Werth bis auf die letzte Einheit richtig ist. In meiner Abhandlung steht \(0{,}24388\), aber dort ist nach der strengen Formel gerechnet, die, wenn man nur fünf Decimalen anwendet, weniger zuverlässig ist als die Näherungsformel. Sehen Sie hier in extenso die Durchführung des vierten Versuchs. Die rothen\HistoricalNoteEight{Die im Original mit rother Tinte geschriebenen Zahlen sind hier unterstrichen.]} Zahlen sind bei allen Versuchen dieselben und die Bedeutung der einzelnen Zahlen steht daneben:
\begingroup
\small
\setlength{\arraycolsep}{5pt}
\renewcommand{\arraystretch}{1.35}
\[
\begin{array}{rrr@{\quad}l}
0{,}24385&&&u\\
\underline{0{,}31365}&&&c\\
\underline{0{,}95443}&&&c''\\
\underline{9{,}74351}&\underline{9{,}91066}&\underline{9{,}22527}&bB,\ b''B'',\ A\\
9{,}74624&0{,}07856&9{,}38712&u+c,\ u+c'',\ u\\
\underline{0{,}17717}&\underline{0{,}05049}&9{,}91567&B\BareRoot{\frac{12}5},\ B''\BareRoot{\frac{12}5},\ \cos\eta\\
9{,}84812&9{,}74970&&\sin\theta,\ \sin\theta''\\
0{,}32905&&&s\\
0{,}30079&&&s''\\
0{,}28713&&&\text{\footnotemark}\\
0{,}61618&&&s+s''
\end{array}
\]
\[
\left.\begin{array}{r}9{,}47145\\0{,}30809\\-22\end{array}\right\}
9{,}77932\ \text{vollkommen genau}
\qquad
\left.\begin{array}{l}k\\\BareRoot{(s+s'')}\\-\BareRoot[5]{\sec Q}\end{array}\right\}\text{Resultat}
\]
\endgroup
\footnotetext{Aus der zweiten Columne meiner Hülfstafel, indem man in die erste mit der Differenz der Logarithmen von \(s\) und \(s''\) eingeht.}

\SourceLeaf{341}\noindent\hspace*{\parindent}.\ .\ .\ .\ .\ .\par
Ich bemerke noch, dass von den oben behuf der ersten Annäherung gemachten Voraussetzungen die erste bei mässigen Zwischenzeiten völlig \textls[150]{legitim} ist. Insofern \(r\), \(r''\) ähnliche Functionen von Grössen sind, die sich nach dem Gesetz der Stetigkeit ändern, kann man näherungsweise das arithmetische Mittel \(\tfrac12(r+r'')\) als dieselbe Function der respectiven Mittel ansehen. Die Ausdehnung auf \(k\) kann freilich in manchen Fällen gar zu gewagt sein, indessen ist der Versuch auf alle Fälle unschädlich. Lässt man die letzte Licenz weg, so wird \(u\) durch folgende Gleichung zu bestimmen sein:
\[
\BareRoot[4]{(bb'')}\cdot\frac{t''-t}{m\BareRoot{18}}=\BareRoot{(uu+AA)}\,.\BareRoot[4]{\left\{\bigl(u+\tfrac12(c+c'')\bigr)^2+bb''BB''\right\}},
\]
die entwickelt auf den 6. Grad steigen würde, aber sich leicht durch Versuche auflösen lässt. Dies wird wohl immer eine wirklich brauchbare Annäherung geben. In unserm Beispiel folgt daraus \(u=0{,}2382\), wie Sie sehen, der Wahrheit schon sehr nahe. Am besten wird es wohl sein, vermittelst dieser Gleichung nur eine oder zwei Ziffern von \(u\) zu berechnen, und dann sofort zu den genauen Versuchen zu schreiten.
\par\medskip\noindent\makebox[\linewidth]{\rule{.18\linewidth}{.35pt}}\par
